\section{Interfaz gráfica en Python}

Se desarrolló una interfaz gráfica en Python para graficar archivos CSV provenientes del osciloscopio y de LTspice. El programa fue diseñado para aceptar archivos temporales y archivos de Bode, detectando automáticamente el tipo de archivo cargado.

\subsection{Funcionalidades implementadas}

La interfaz permite realizar las siguientes operaciones:

\begin{itemize}[itemsep=0.2em]
    \item cargar uno o varios archivos CSV;
    \item arrastrar y soltar archivos sobre la ventana;
    \item detectar si el archivo corresponde a una señal temporal o a un diagrama de Bode;
    \item impedir la mezcla accidental de datos temporales y datos de Bode en un mismo gráfico;
    \item graficar múltiples canales, incluso si provienen de archivos distintos;
    \item modificar escala vertical, offset vertical y offset horizontal de cada canal;
    \item alinear señales temporales a $t=0$ para comparar mediciones con simulaciones;
    \item desplazar un canal en el eje horizontal mediante teclado, mouse o ingreso numérico;
    \item agregar cursores sobre una curva, con cálculo automático de la coordenada $Y$;
    \item comparar cursores mediante diferencias $\Delta X$ y $\Delta Y$;
    \item activar modo XY para obtener figuras de Lissajous;
    \item marcar máximos y mínimos;
    \item exportar el gráfico final a PDF.
\end{itemize}

\subsection{Organización del programa}

El archivo principal de la interfaz se incluye en la carpeta \texttt{codigo/}. Se utilizaron las bibliotecas \texttt{pandas}, \texttt{numpy}, \texttt{matplotlib} y \texttt{tkinter}. De forma opcional, se empleó \texttt{tkinterdnd2} para habilitar la carga de archivos mediante arrastrar y soltar.

\begin{lstlisting}[style=pythonstyle, caption={Comando de instalación de dependencias.}]
py -m pip install pandas numpy matplotlib tkinterdnd2
\end{lstlisting}

\begin{lstlisting}[style=pythonstyle, caption={Ejecución de la interfaz.}]
py gui_tc1_user_friendly_pdf.py
\end{lstlisting}

\subsection{Modo temporal y modo Bode}

Una mejora importante fue separar el funcionamiento en dos modos excluyentes. Si la interfaz se encuentra en modo Bode y se intenta cargar una señal temporal, el programa avisa al usuario y pregunta si desea borrar los Bode cargados y cambiar a modo temporal. Análogamente, si se está trabajando en modo temporal y se carga un Bode, el programa solicita confirmación antes de limpiar los datos actuales. Esto evita superponer datos con ejes físicos incompatibles.

En modo temporal, el eje horizontal representa tiempo y puede expresarse en segundos, milisegundos, microsegundos o nanosegundos. En modo Bode, el eje horizontal representa frecuencia en hertz y se grafica automáticamente en escala logarítmica, separando magnitud y fase en dos subgráficos.
